\section{BDD-сценарії}
\label{sec:bdd}

Нижче наведено BDD-сценарії у форматі Gherkin для ключових функцій системи.

\begin{codelisting}
\caption{BDD-сценарії системи}
\label{lst:bdd}
\begin{skedcode}
Сценарій: Реєстрація нового користувача
  Given відкрита сторінка реєстрації
  And email "test@example.com" ще не зареєстрований у системі
  When користувач вводить email "test@example.com", пароль та нікнейм
    і підтверджує реєстрацію
  Then користувач отримує доступ до порожнього особистого каталогу
  And повторна спроба реєстрації з тим самим email показує повідомлення
    "Акаунт із цим email вже існує"
\end{skedcode}
\end{codelisting}
\sked{K7}

\begin{codelisting}
\caption{Вхід до акаунту}
\label{lst:bdd-login}
\begin{skedcode}
Сценарій: Вхід до акаунту з коректними даними
  Given користувач зареєстрований у системі
  When він вводить правильний email і пароль та натискає "Увійти"
  Then він перенаправляється до свого особистого каталогу
  And бачить свої раніше додані елементи
\end{skedcode}
\end{codelisting}
\sked{K8}

\begin{codelisting}
\caption{Повторний перегляд з іншою оцінкою}
\label{lst:bdd-session}
\begin{skedcode}
Сценарій: Повторний перегляд фільму з іншою оцінкою
  Given елемент контенту "Фільм X" вже має одну Session з оцінкою 9
  When користувач додає нову Session для того самого елемента
    з оцінкою 5 і нотаткою "на другий раз не сподобалось"
  Then елемент контенту має дві Sessions, кожна зі своєю датою,
    оцінкою та нотаткою
  And обидві оцінки зберігаються окремо, жодна не перезаписується
\end{skedcode}
\end{codelisting}
\sked{K18}

\begin{codelisting}
\caption{Елемент у кількох добірках}
\label{lst:bdd-collection}
\begin{skedcode}
Сценарій: Елемент належить кільком добіркам одночасно
  Given користувач має елемент "Книга Y" у каталозі
  And існують добірки "Улюблене" та "2026 рік"
  When користувач додає "Книга Y" до обох добірок
  Then "Книга Y" відображається у вмісті обох добірок
  And видалення "Книга Y" з добірки "2026 рік" не видаляє її
    з "Улюблене" чи з каталогу
\end{skedcode}
\end{codelisting}
\sked{K21,K25}

\begin{codelisting}
\caption{Автозаповнення дати завершення}
\label{lst:bdd-date}
\begin{skedcode}
Сценарій: Автозаповнення дати завершення
  Given користувач редагує елемент зі статусом "у процесі"
  When користувач змінює статус на "завершено"
    і не вказує дату вручну
  Then система автоматично встановлює дату завершення рівною
    поточній даті
\end{skedcode}
\end{codelisting}
\sked{K27}

\begin{codelisting}
\caption{Фільтрування за тегом}
\label{lst:bdd-tag}
\begin{skedcode}
Сценарій: Фільтрування каталогу за тегом
  Given у каталозі є три елементи: два з тегом "фентезі" і один без
  When користувач обирає фільтр за тегом "фентезі"
  Then відображаються лише два елементи, позначені цим тегом
  And при скиданні фільтра знову відображаються всі три елементи
\end{skedcode}
\end{codelisting}
\sked{K17}
